\scanpage{215}
(Résulte de la conjonction des conj de Hodge et de Tate \inserted{\tikz[baseline=(a.base)]\node[draw,ellipse,inner sep=1pt](a){sous forme faible.};} \unsure{?}, du moins si $X$ \uline{\unsure{normal|connexe|réduit}}). \uline{\uuline{\deleted{\unsure{Caractérisation}}Caractériser \deleted{de} l'image}} ??
\margin{$\vert\,\vert\,\vert\,\vert\,\vert$}

M conclusion si $X$ défini sur $k\subset\mathbb C$, $k$ alg. clos, p. ex. $k=\overline{\mathbb Q}$.

\begin{center}\uline{Analyser l'exemple de Serre de} $\pi_1$ \uline{non isomorphes}\end{center}

[\uline{Cas corps de base} $\mathbb R$ : tout pt \deleted{réel} \inserted{fermé $a$} définit une $\unsure{\xi|\zeta}_a\in T_{\unsure{\alpha|a}}(\mathbb Z)$. Tous ces $\xi_a$ sont isomorphes si \deleted{\unsure{a}} $X$ est \deleted{\unsure{géom.} \uline{connexe}, \unsure{sinon}} \inserted{\unsure{irrid|irréd} (pas néc. géom \unsure{irrid|irréd})} \deleted{il peut y en avoir deux non isom.}

Si $a$ est un pt réel, alors $\xi_a$ admet un \unsure{autom} $\mathrm{frob}_{\infty}$ d'ordre $2$, qui avec la comp. conn. de $G(X,\xi)$ engendre $G(X,\xi)$.
\margin{\diagram{\begin{tikzcd}[ampersand replacement=\&,row sep=.5em,column sep=.6em]
{} \& X\arrow[dd,dash]\\
\mathbb C\arrow[ur,dash]\arrow[dr,dash] \& {}\\
{} \& \mathbb R
\end{tikzcd}}}
Le foncteur
\[
\begin{aligned}
M\ \unsure{\sim|\leadsto}\ (\unsure{\xi(M_{\xi})|\xi(M)_{\xi}|\xi(M_{\mathbb C})},\ \unsure{+^{\prime}|+}&\text{structure de Hodge}+\text{\unsure{réellification|réalification} définie}\\
&\text{par }\mathrm{frob}_{\infty}\,)
\end{aligned}
\]
est pleinement fidèle. ]

\margin{$\vert\,\vert\,\vert$}
\uline{Pb} étudier comment varie la structure de Hodge pour un pt \inserted{complexe} variable, \uline{en fonction de la donnée \unsure{départ|discret|discrète|définie}} en le pt $a_0$,
% archive sheet number: 214

